\section{Conclusion}\label{sec:conclusion}

This paper reframes meibomian gland segmentation from ``design another segmentation network'' to ``build comparable, auditable evidence''. On the public MGD-1K dataset we establish a patient-level split (\nTrain/\nVal/\nTest{} images; \nTrainPatients/\nValPatients/\nTestPatients{} patients) and record both a file-list signature and an assignment signature, so that ``same filenames, different assignment'' is no longer an untraceable hidden variable.

Methodologically, we provide three reusable findings. First, leakage is quantifiable: the official \officialTestImages-case Meiboscore test list comes from \officialTestPatients{} patients, and \officialLeakImages{} of the remaining 800 images come from the same patients; even a patient-level training set still contains \officialLeakPatients/\officialTestPatients{} test patients. Second, tiling is an experimental variable on par with the backbone: with identical weights, changing only tile height moves Dice from 0.370 to 0.655; MTE averages probabilities from 256/384/512 tiles, but in this run it mainly improved HD95 (68.05→61.85) rather than Dice, so it is a boundary correction rather than an accuracy gain. Multi-tile variance reaches only 0.639 AUROC for failure detection, and TTA8's 1.000 is affected by label--signal coupling; neither is a standalone triage rule. Third, pixel metrics and clinical endpoints must be reported together: the QWK between segmentation-derived Meiboscore and expert grading is only 0.131, so the current R\_loss→Meiboscore chain cannot replace expert grading; upper and lower eyelids must also be reported separately.

Architecturally, MGD-Former keeps the 1/4 feature from a MiT-B2 encoder, uses a fusion decoder at stride 4, and supervises three scales; ADE20K pretraining targets the cold start on \nTotal{} cases. It is one of five routes, not the sole selling point. We report no unmeasured accuracy claim; all values have been injected from the unified training and patient-level independent test results, supported by confidence intervals, paired tests, and multiple metrics (Dice, clDice, HD95, ASSD, clinical QWK).
